\documentclass[11pt]{article}
\usepackage{auditstyle}
\usepackage{caption}
\fancyhead[L]{Part II source-evidence dossier}
\fancyhead[R]{9 August 2026}
\hypersetup{pdftitle={Part II citation screenshot dossier},pdfauthor={Proof audit}}
\newcommand{\shot}[1]{%
  \begin{center}
  \includegraphics[width=.88\textwidth,height=.73\textheight,keepaspectratio]{#1}
  \end{center}}
\newcommand{\evidence}[4]{%
  \section*{#1}
  \textbf{Source location.} #2

  \textbf{What the page establishes.} #3

  \textbf{What it does not establish.} #4
}

\begin{document}

\begin{center}
{\Large\bfseries Part II Citation Screenshot Dossier}\\[4pt]
{\large Wavelet resolution / operator-algebraic renormalization}\\[10pt]
\textbf{Companion to the Lamport proof audit}\\[10pt]
9 August 2026
\end{center}

\begin{resultbox}
This dossier records the exact pages used for the source-scope and normalization assertions
in the Part II Lamport proof package.  A page image is
evidence only for the text visible on that page.  Every inference below distinguishes a
positive statement in the source from a stronger statement that the source does not make.
The mathematical no-go theorem in the companion document is proved independently and
does not rest on a numerical experiment or on absence-of-theorem reasoning.
\end{resultbox}

\subsection*{Source manifest}

\begin{enumerate}
\item V. Morinelli, G. Morsella, A. Stottmeister, Y. Tanimoto,
\emph{Scaling limits of lattice quantum fields by wavelets}, CMP 387 (2021),
1299--1367.  Local file SHA-256:
{\footnotesize\ttfamily 3882444339413145bae290b8213582cb\newline
9a388c79b16ac2e2ee4e15dcac79c01e}.
\href{https://doi.org/10.1007/s00220-021-04152-5}{Publisher DOI}.

\item G. Battle, P. Federbush,
\emph{Ondelettes and phase cell cluster expansions, a vindication}, CMP 109
(1987), 417--419.  Local file SHA-256:
{\footnotesize\ttfamily d68ceaf1e27ca82f0db995e94a7168eb\newline
2aa448dddbe82d83a553bf0395b94c07}.
\href{https://doi.org/10.1007/BF01206144}{Publisher DOI}.

\item A. Stottmeister, V. Morinelli, G. Morsella, Y. Tanimoto,
\emph{Operator-algebraic renormalization and wavelets}, arXiv:2002.01442v2.
Local file SHA-256:
{\footnotesize\ttfamily 2db5b30468cba5f073769774838a8925\newline
d36060809b65ecaf9e57cb5b753c4ec2}.
\href{https://arxiv.org/abs/2002.01442}{arXiv record}.

\item B. Nachtergaele, H. Raz, B. Schlein, R. Sims,
\emph{Lieb--Robinson bounds for harmonic and anharmonic lattice systems},
CMP 286 (2009), 1073--1098.  Local file SHA-256:
{\footnotesize\ttfamily 613ff5cc8af3f7b9734a2ca1912f33624\newline
b59050204e108071c4d200285179114}.
\href{https://arxiv.org/abs/0712.3820}{arXiv record}.
\end{enumerate}

\textbf{Integrity check.}  Each screenshot was rendered directly from the corresponding
local PDF.  No OCR text was substituted for a page image.  PDF page numbers below are
one-based file page numbers; where visible, the printed journal page is also reported.

\clearpage

\evidence{E1. The OAR limit is projective at fixed coarse scale}
{MMST, PDF page 7, equations (2.8)--(2.11).}
{Observable algebras form an inductive limit.  A finer state is pulled back to each
fixed coarse algebra; the horizontal limit is then required at that fixed scale, and a
projectively consistent family defines the limiting state.}
{This page states neither convergence in the norm of local state spaces nor a bound
uniform over a coarse scale that is allowed to equal the fine scale.}
\shot{evidence/mmst_pdf_page_007.png}

\clearpage

\evidence{E2. The continuum map is an observable embedding}
{MMST, PDF page 25, Proposition 4.1.}
{The symplectic one-particle map induces an injective $*$-homomorphism
$\beta_L:\mathcal W_{\infty,L}\to\mathcal W(h_L)$ and gives its action on Weyl
generators.}
{An injective homomorphism of observable algebras does not canonically push a state in
the reverse direction.  This proposition does not put a lattice Hamiltonian and a
continuum Hamiltonian on one Hilbert space or one common form domain.}
\shot{evidence/mmst_pdf_page_025.png}

\clearpage

\evidence{E3. Identification is pointwise on embedded observables}
{MMST, PDF page 28, Lemma 4.3 and the paragraph following it.}
{For every observable $W$ belonging to a fixed finite-scale Weyl algebra, the limiting
state equals the continuum vacuum evaluated on $\beta_L(W)$.  The text then invokes
projective consistency to obtain a state on the inductive limit.}
{The equality identifies the already-taken horizontal limit.  It is not a quantitative
diagonal estimate comparing a level-$N$ state on all level-$N$ local observables with a
continuum state.}
\shot{evidence/mmst_pdf_page_028.png}

\clearpage

\evidence{E4. Proposition 4.5 states the free scaling limit}
{MMST, PDF page 29, Proposition 4.5.}
{For the specified Daubechies scaling maps and mass renormalization, the sequence of
free lattice states has a scaling limit, and that limit agrees with the continuum vacuum
on observables coming from some finite scale.}
{No rate $C\varepsilon_N^\theta$ and no supremum over the finest-scale unit ball appears.
The proposition is about free ground states, not interacting $\phi^4_2$ ground states.}
\shot{evidence/mmst_pdf_page_029.png}

\clearpage

\evidence{E5. The free continuum time-zero net is recovered}
{MMST, PDF page 35, Theorem 4.11 and Remarks 4.12--4.13.}
{The scaling-limit state yields a continuum time-zero net unitarily equivalent to the
massive free-field net.  Remark 4.12 separately describes a diagonal sequence used for
the dynamics.}
{The theorem does not identify the finite-scale and continuum Hamiltonians as operators
on a common Hilbert space, and it does not prove norm convergence of diagonal states.}
\shot{evidence/mmst_pdf_page_035.png}

\clearpage

\evidence{E6. The cited Lieb--Robinson estimate is harmonic}
{MMST, PDF page 36, Lemma 4.14.}
{A commutator estimate is quoted and rescaled for the free harmonic lattice Hamiltonian
and lattice Weyl observables.  Its constants and distance dependence are displayed.}
{The lemma is not an interacting quartic Lieb--Robinson theorem.  A bound on the
derivative of the one-particle dispersion alone cannot replace the commutator estimate.}
\shot{evidence/mmst_pdf_page_036.png}

\clearpage

\evidence{E7. The LR speed bound is not the exact group velocity}
{MMST, PDF page 38, Remark 4.17.}
{The authors explicitly distinguish the scaling-limit Lieb--Robinson velocity obtained
from their estimate from the physical propagation speed; the displayed optimum is
approximately $3.59$, whereas the physical speed is $1$.}
{Consequently the exact free group-velocity calculation in the companion proof is a
useful check, but it is not itself the LR theorem needed for an interacting model.}
\shot{evidence/mmst_pdf_page_038.png}

\clearpage

\evidence{E8. Interacting wavelet scaling is stated as future work}
{MMST, PDF page 54, Section 6.1.}
{The text says that recovery of interacting dynamics should be generalized to the
wavelet renormalization group and that one should work with scaling limits of interacting
lattice ground states.}
{The modal language is prospective.  This page does not prove an interacting wavelet
ground-state limit, an $N$-uniform cluster expansion, an $N$-uniform physical gap, or the
mixed-vertex estimate required by the original Part II program.}
\shot{evidence/mmst_pdf_page_054.png}

\clearpage

\evidence{E9. Battle--Federbush describes mechanism and an earlier error}
{Battle--Federbush (1987), PDF page 1 / journal page 417.}
{The abstract presents wavelet phase-cell expansions as a way to improve earlier work
and explicitly says that an early error is excised.  Section 1 begins with Meyer-type
orthonormal modes.}
{This introductory page contains no theorem in the MMST OAR lattice variables and no
scale-uniform constants CE1$_N$ or CE2$_N$.}
\shot{evidence/bf87_pdf_page_001.png}

\clearpage

\evidence{E10. The wavelet bases differ from the MMST finite-scale action}
{Battle--Federbush (1987), PDF page 2 / journal page 418, Sections 1--5.}
{The page distinguishes Meyer functions, sharply localized nonsmooth functions, and
exponentially localized bases; it records their localization and moment properties and
starts the phase-cell field expansion.}
{These basis properties do not identify the phase-cell action with the nearest-neighbour
MMST lattice Hamiltonian.  They also do not provide the OAR Wick constants, local
observable norms, or a mixed lattice--continuum interpolation.}
\shot{evidence/bf87_pdf_page_002.png}

\clearpage

\evidence{E11. The phase-cell free action is diagonal in different variables}
{Battle--Federbush (1987), PDF page 3 / journal page 419, Section 5.}
{The phase-cell modes are chosen as
$u_k=(-\Delta+M^2)^{-1/2}\psi_k$, and the authors emphasize that the free action is
diagonal in those variables.  The note explains how replacing the earlier basis repairs
the cited phase-cell constructions.}
{The page does not state the two uniform estimates demanded in Part II for a
nearest-neighbour OAR action, nor does it prove a uniform interacting Hamiltonian gap or
an OAR local-state convergence rate.}
\shot{evidence/bf87_pdf_page_003.png}

\clearpage

\evidence{E12. The short OAR article also labels the interacting step as prospective}
{Stottmeister--Morinelli--Morsella--Tanimoto, arXiv:2002.01442v2, PDF page 5.}
{The conclusion says the framework can include interacting lattice systems but would
need analytic or numerical expansions; it describes an extension to the wavelet setting
as expected and asks whether explicit error bounds can be proved.}
{This is an outlook statement, not a proof of the interacting wavelet scaling limit or
of uniform cluster-expansion, gap, Lieb--Robinson, or local-state-norm estimates.}
\shot{evidence/smmt_pdf_page_005.png}

\clearpage

\evidence{E13. MMST fixes the lattice Weyl and Fourier normalizations}
{MMST, PDF page 13, equations (3.5)--(3.10).}
{In dimension \(d=1\), equation (3.7) reads
\(\xi=\varepsilon_Nq+ip\); equations (3.8)--(3.9) give the symplectic form
\(\varepsilon_N\sum_x(qp'-pq')\) and the normalized discrete Fourier variables.
Equation (3.10) states how a symplectic one-particle map induces the observable
embedding.}
{This page does not choose a particular low-pass mask or define an interacting
Hamiltonian.  Those data must be supplied separately.}
\shot{evidence/mmst_pdf_page_013.png}

\clearpage

\evidence{E14. The wavelet refinement map and its symplectic property are explicit}
{MMST, PDF page 16, Definition 3.3 and Lemma 3.4.}
{For \(d=1\), equation (3.25) gives the one-step factor \(\sqrt2\), the translated
low-pass coefficients \(h_n\), periodization, and composition across scales.
Lemma 3.4 states that these maps preserve the scale-dependent symplectic forms.}
{The lemma concerns observable kinematics.  It does not say that the
nearest-neighbour Hamiltonians are mapped into one another or that interacting ground
states converge.}
\shot{evidence/mmst_pdf_page_016.png}

\clearpage

\evidence{E15. The \(D4\) scaling function has the support used in the model sheet}
{MMST, PDF page 20, Section 3.3.4.}
{For the Daubechies member \(K=2\), the source states the filter cutoff \(n\ge4\),
the scaling-function support \([0,3]\), and regularity sufficient for the continuum
field and momentum construction.}
{This page does not display the four numerical coefficients.  The model sheet fixes
them explicitly and verifies the sum rule and QMF identity by direct algebra rather
than importing them without a check.}
\shot{evidence/mmst_pdf_page_020.png}

\clearpage

\evidence{E16. The nearest-neighbour free Hamiltonian and vacuum covariance are fixed}
{MMST, PDF page 21, equations (4.1)--(4.5).}
{The page displays the nearest-neighbour lattice Hamiltonian, its dispersion, the Weyl
vacuum characteristic functional, and the separate \(q\)- and \(p\)-covariance
multipliers.  In \(d=1\) and after the stated physical-mass reparametrization this gives
\(\gamma_N(k)^2=m^2+4\varepsilon_N^{-2}\sin^2(\varepsilon_Nk/2)\).}
{These are free formulas.  They neither fix a Wick-quartic convention nor prove an
interacting resolution limit.}
\shot{evidence/mmst_pdf_page_021.png}

\clearpage

\evidence{E17. The fixed-coarse free covariance contains the finite mask product}
{MMST, PDF page 26, equations (4.17)--(4.20) and Lemma 4.2.}
{Equation (4.17) displays the finite product
\(\prod_{n=1}^M m_0(\varepsilon_{N+n}k)\) in the pulled-back free covariance, and
equation (4.19) displays its scaling-function limit.  Lemma 4.2 proves weak-star
convergence at each fixed coarse scale.}
{The source proves convergence, not the explicit power rate.  The companion rate
document derives that rate from the exact \(D4\) factorization and gives every
constant.}
\shot{evidence/mmst_pdf_page_026.png}

\clearpage

\evidence{E18. The cited anharmonic Lieb--Robinson theorem excludes the Wick quartic}
{Nachtergaele--Raz--Schlein--Sims, PDF page 19, Theorem 4.1, equations (4.4)--(4.5).}
{The theorem explicitly assumes \(V\in C^1(\mathbb R)\),
\(V'\in L^1(\mathbb R)\), and
\(\int |w|\,|\widehat{V'}(w)|\,dw<\infty\), before stating the commutator bound.}
{For \(V_c(u)=\lambda(u^4-6cu^2+3c^2)\),
\(|V_c'(u)|\ge2\lambda|u|^3\) on \(|u|\ge\sqrt{6c}\), so the displayed
\(L^1\) hypothesis fails.  This excludes this theorem as a citation route; it does not
prove that no quartic Lieb--Robinson estimate can exist.}
\shot{evidence/nrss_pdf_page_019.png}

\clearpage

\section*{Addendum (12 August 2026): the V3 import cluster}

Phase V3 (\texttt{v3\_continuum\_package.tex}) resolves Statement A1's
``import audit or proof'' deliverable as \emph{proof}, on the scoping
note's fallback route; the H{\o}egh-Krohn (1974) and
Figari--H{\o}egh-Krohn--Nappi (1975) candidates are \emph{not} consumed.
The only printed import is the $\Gamma(A)$ cluster of Simon, \emph{The
$P(\phi)_2$ Euclidean (Quantum) Field Theory} (PDF in \texttt{refs/}),
the same source and theorem (I.17) as Part I's Assumption
\texttt{A:hyper}. Evidence pages follow.

\evidence{E19. $\Gamma(A)$ is positivity preserving (Theorem I.12)}
{Simon book, PDF page 48 (printed page 26).}
{For any contraction $A$ of real Hilbert spaces, $\Gamma(A)$ maps
nonnegative functions to nonnegative functions.}
{Nothing about strict positivity (that is Theorem I.16) or about
$L^p$ bounds (Corollary I.15, Theorem I.17).}
\shot{dossier_images/simon-gamma-p48-048.png}

\clearpage

\evidence{E20. $\Gamma(A)$ is an $L^p$ contraction (Corollary I.15)}
{Simon book, PDF page 53 (printed page 31).}
{For a contraction $A$ on a real Hilbert space, $\Gamma(A)$ is a
contraction on each $L^p(M,d\mu)$.}
{No hypercontractive improvement (that is Theorem I.17).}
\shot{dossier_images/simon-gamma-p53-053.png}

\clearpage

\evidence{E21. $\Gamma(A)$ is positivity improving for $\|A\|<1$
(Theorem I.16)}
{Simon book, PDF page 55 (printed page 33); proof concludes on PDF page
56.}
{For a strict contraction $A$, $\Gamma(A)$ is positivity improving.
Used in V3, Theorem 7.1 for the two-time pair positivity
$\langle\mathbf 1_A,\Gamma(e^{-t\gamma})\mathbf 1_B\rangle>0$.}
{False at $\|A\|=1$ (the source's Remark 3); V3 applies it only with
$\|e^{-t\gamma}\|=e^{-tm}<1$.}
\shot{dossier_images/simon-gamma-p55-055.png}

\clearpage

\evidence{E22. Nelson's hypercontractive bound (Theorem I.17)}
{Simon book, PDF page 56 (printed page 34), display (I.42)--(I.43);
remarks continue on PDF page 57.}
{For a contraction $A:\mathcal H_1\to\mathcal H_2$ with
$\|A\|\le(p-1)^{1/2}(q-1)^{-1/2}$, $1<p\le q<\infty$:
$\|\Gamma(A)\psi\|_q\le\|\psi\|_p$. This is the abstract form quoted by
Part I's \texttt{A:hyper} and consumed by V3's Lemma 5.1 (ladder) and
the chaos moment bounds.}
{The optimality remark (violation of (I.42) destroys boundedness) is
not used.}
\shot{dossier_images/simon-gamma-p56-056.png}

\clearpage

\evidence{E23. Theorem I.17, continuation page}
{Simon book, PDF page 57 (printed page 35), Remarks 3--4 and the
reduction $A=B\|A\|$ to the diagonal case.}
{The reduction of the general contraction to $A=e^{-t}$ via
$\Gamma(A)=\Gamma(B)\Gamma(\|A\|)$ with $\Gamma(B)$ an $L^q$
contraction --- the composition pattern V3 uses for
$\Gamma(e^{-\tau\gamma})=\Gamma(e^{-\tau(\gamma-m)})
\Gamma(e^{-\tau m})$.}
{The weakened-result proof route sketched there is not consumed; only
the stated Theorem I.17 is.}
\shot{dossier_images/simon-gamma-p57-057.png}

\end{document}
